\documentclass[a4paper,12pt]{article}

\usepackage[T2A]{fontenc}
\usepackage[utf8]{inputenc}
\usepackage[russian]{babel}

\usepackage{amsmath,amssymb,mathtools}
\usepackage{paratype}
\renewcommand{\familydefault}{\sfdefault}

\usepackage{icomma}
\usepackage[margin=20mm]{geometry}
\usepackage[nopatch=footnote]{microtype}
\usepackage{tikz}
\usetikzlibrary{calc,arrows.meta,positioning,shapes.geometric}

\usepackage{tabularx,booktabs}
\usepackage{enumitem}
\usepackage{hyperref}
\usepackage{parskip}
\usepackage{fancyhdr}
\usepackage{setspace}
\usepackage{xcolor}

\definecolor{accent}{HTML}{008080}
\definecolor{darkaccent}{HTML}{004D4D}
\definecolor{textgray}{HTML}{333333}

\hypersetup{hidelinks}

\onehalfspacing
\setlength{\parindent}{0pt}
\setlength{\headheight}{25pt}

\setlist[itemize]{
    leftmargin=1.7em,
    itemsep=0.35em,
    topsep=0.35em
}

\setlist[enumerate]{
    leftmargin=1.9em,
    itemsep=0.5em,
    topsep=0.4em
}

\pagestyle{fancy}
\fancyhf{}
\fancyhead[L]{\small\color{textgray}Средняя скорость}
\fancyhead[R]{\small\color{textgray}Кирилл Зиновьев, 7 класс}
\fancyfoot[C]{\small\thepage}

\renewcommand{\headrulewidth}{0.3pt}
\renewcommand{\headrule}{
    \hbox to\headwidth{
        \color{accent}
        \leaders\hrule height \headrulewidth\hfill
    }
}

\newcommand{\avgv}{v_{\text{ср}}}
\newcommand{\unit}[1]{\,\text{#1}}

\tikzset{
    flowstart/.style={
        ellipse,
        draw=darkaccent,
        line width=0.8pt,
        minimum width=38mm,
        minimum height=10mm,
        align=center,
        inner sep=4pt
    },
    flowaction/.style={
        rounded corners=2mm,
        draw=darkaccent,
        line width=0.8pt,
        minimum width=70mm,
        text width=68mm,
        minimum height=12mm,
        align=center,
        inner sep=5pt
    },
    flowdecision/.style={
        diamond,
        aspect=2.6,
        draw=darkaccent,
        line width=0.8pt,
        text width=40mm,
        align=center,
        inner sep=2pt
    },
    flowarrow/.style={
        -{Stealth[length=3mm,width=2mm]},
        line width=0.8pt,
        draw=darkaccent
    },
    flowlabel/.style={
        font=\small,
        fill=white,
        inner sep=1pt,
        text=darkaccent
    }
}

\begin{document}

% =========================================================
% ТИТУЛЬНЫЙ ЛИСТ
% =========================================================

\begin{titlepage}
\thispagestyle{empty}
\centering

\vspace*{0.14\textheight}

{\Large\color{darkaccent}\bfseries Чек-лист\par}

\vspace{0.7em}

{\huge\bfseries
Средняя скорость: путь, время и перевод в СИ
\par}

\vspace{1.4em}

{\large 7 класс\par}

\vspace{2.2em}

{\large Лёвин Артём Александрович\par}

\vspace{0.4em}

{\large эксклюзивно для Кирилла Зиновьева\par}

\vspace{1.4em}

{\large \today\par}

\vfill

\begin{tikzpicture}[remember picture,overlay]
\draw[
    accent,
    line width=1.1pt
]
([xshift=18mm,yshift=20mm]current page.south west)
.. controls
([xshift=55mm,yshift=31mm]current page.south west)
and
([xshift=-55mm,yshift=13mm]current page.south east)
..
([xshift=-18mm,yshift=24mm]current page.south east);
\end{tikzpicture}

\end{titlepage}

% =========================================================
% ОСНОВНОЙ ЧЕК-ЛИСТ
% =========================================================

\section*{1. Что нужно знать}

\subsection*{Средняя путевая скорость}

Если тело проходит несколько участков пути, средняя путевая скорость равна отношению
\textbf{всего пройденного пути} ко \textbf{всему затраченному времени}:

\[
\boxed{
\avgv=
\frac{s_{\text{общ}}}{t_{\text{общ}}}
}
\]

Для трёх участков:

\[
\boxed{
\avgv=
\frac{s_1+s_2+s_3}
     {t_1+t_2+t_3}
}
\]

Индексы $1$, $2$, $3$ помогают отделять данные разных участков движения:

\[
s_1,\ t_1
\]

относятся к первому участку,

\[
s_2,\ t_2
\]

--- ко второму и т.~д.

Индекс записывается справа внизу у обозначения величины.

\subsection*{Система СИ}

Перед подстановкой чисел данные удобно привести к единицам СИ.

Для задач занятия нужны три основных перевода:

\[
1\unit{км}=1000\unit{м},
\]

\[
1\unit{мин}=60\unit{с},
\]

\[
1\unit{ч}=3600\unit{с}.
\]

Следовательно,

\[
2{,}5\unit{км}
=
2{,}5\cdot1000
=
2500\unit{м},
\]

\[
10\unit{мин}
=
10\cdot60
=
600\unit{с}.
\]

Если путь уже дан в метрах, а время --- в секундах, дополнительный перевод выполнять не требуется.

\subsection*{Пример 1. Три участка пути}

Сноубордист прошёл:

\[
\begin{aligned}
s_1&=2{,}5\unit{км},
&
t_1&=10\unit{мин},
\\
s_2&=800\unit{м},
&
t_2&=1\unit{мин},
\\
s_3&=1{,}2\unit{км},
&
t_3&=4\unit{мин}.
\end{aligned}
\]

Переведём данные в СИ:

\[
s_1=2500\unit{м},
\qquad
s_2=800\unit{м},
\qquad
s_3=1200\unit{м},
\]

\[
t_1=600\unit{с},
\qquad
t_2=60\unit{с},
\qquad
t_3=240\unit{с}.
\]

Используем формулу:

\[
\avgv=
\frac{s_1+s_2+s_3}
     {t_1+t_2+t_3}.
\]

Подставляем данные:

\[
\avgv=
\frac{2500+800+1200}
     {600+60+240}.
\]

Складываем отдельно путь и время:

\[
s_{\text{общ}}
=
2500+800+1200
=
4500\unit{м},
\]

\[
t_{\text{общ}}
=
600+60+240
=
900\unit{с}.
\]

Получаем:

\[
\avgv=
\frac{4500}{900}
=
5\unit{м/с}.
\]

\[
\boxed{\avgv=5\unit{м/с}}
\]

\subsection*{Что важно понимать}

\begin{itemize}

\item
В формуле для средней путевой скорости учитывается суммарная длина всех пройденных участков.

\item
Уклон трассы сам по себе формулу не меняет. Если спортсмен прошёл $800$ метров, в суммарный путь добавляются именно эти $800$ метров.

\item
Если тело развернулось и часть пути прошло обратно, для \textbf{средней путевой скорости} длины всех участков по-прежнему складываются.

\item
Для средней скорости как векторной величины используется перемещение. Это отдельное понятие, которое важно отличать от средней путевой скорости.

\item
Обычное среднее арифметическое скоростей отдельных участков в общем случае применять нельзя.

\end{itemize}

\subsection*{Пример 2. Перевод единиц уже не нужен}

Пусть

\[
s_1=10\unit{м},
\qquad
s_2=15\unit{м},
\qquad
s_3=5\unit{м},
\]

\[
t_1=60\unit{с},
\qquad
t_2=40\unit{с},
\qquad
t_3=50\unit{с}.
\]

Все величины уже представлены в метрах и секундах.

Поэтому сразу используем формулу:

\[
\avgv=
\frac{s_1+s_2+s_3}
     {t_1+t_2+t_3}.
\]

\[
\avgv=
\frac{10+15+5}
     {60+40+50}
=
\frac{30}{150}
=
0{,}2\unit{м/с}.
\]

\[
\boxed{\avgv=0{,}2\unit{м/с}}
\]

\subsection*{Пример 3. Часы, минуты и километры}

Пусть

\[
s_1=1{,}5\unit{км},
\qquad
s_2=2\unit{км},
\qquad
s_3=3\unit{км},
\]

\[
t_1=1\unit{ч},
\qquad
t_2=50\unit{мин},
\qquad
t_3=40\unit{мин}.
\]

Переводим путь:

\[
s_1=1500\unit{м},
\qquad
s_2=2000\unit{м},
\qquad
s_3=3000\unit{м}.
\]

Переводим время:

\[
t_1=3600\unit{с},
\qquad
t_2=3000\unit{с},
\qquad
t_3=2400\unit{с}.
\]

Тогда

\[
\avgv=
\frac{1500+2000+3000}
     {3600+3000+2400}.
\]

\[
\avgv=
\frac{6500}{9000}
=
\frac{13}{18}\unit{м/с}
\approx
0{,}72\unit{м/с}.
\]

\[
\boxed{
\avgv\approx0{,}72\unit{м/с}
}
\]

\subsection*{Оформление решения}

Полезно придерживаться одного порядка записи:

\begin{enumerate}

\item
Записать \textbf{Дано}.

\item
Разделить данные разных участков индексами:
$s_1$, $s_2$, $s_3$, $t_1$, $t_2$, $t_3$.

\item
Записать искомую величину:
$\avgv$.

\item
Проверить единицы измерения и выполнить необходимые переводы в СИ.

\item
Записать общую формулу.

\item
Только после формулы подставить числовые значения.

\item
Выполнить вычисления.

\item
Записать \textbf{Ответ} вместе с единицей измерения.

\end{enumerate}

Если Вы начали указывать единицы измерения в промежуточных вычислениях, важно делать это последовательно.

При исправлении записи достаточно аккуратно зачеркнуть ошибочное значение одной линией и рядом написать правильное.

% =========================================================
% БЛОК-СХЕМА
% =========================================================

\clearpage
\thispagestyle{empty}

\begin{center}
{\Large\bfseries\color{darkaccent}
Алгоритм: средняя скорость на нескольких участках}
\end{center}

\vspace{8mm}

\begin{center}

\begin{tikzpicture}[node distance=10mm and 18mm]

\node[flowstart] (start)
{Начало};

\node[flowaction,below=of start] (data)
{Запишите $s_1,s_2,\ldots$ и $t_1,t_2,\ldots$;
укажите искомую $\avgv$};

\node[flowdecision,below=13mm of data] (si)
{Все данные уже в СИ?};

\node[flowaction,below=16mm of si] (sum)
{Найдите
$s_{\text{общ}}=\sum s_i$
и
$t_{\text{общ}}=\sum t_i$};

\node[flowaction,left=22mm of sum] (convert)
{Переведите путь в метры,
время в секунды};

\node[flowaction,below=of sum] (formula)
{Запишите формулу
$\displaystyle
\avgv=
\frac{s_{\text{общ}}}{t_{\text{общ}}}$};

\node[flowaction,below=of formula] (calc)
{Подставьте числа и выполните вычисления};

\node[flowaction,below=of calc] (answer)
{Запишите ответ в $\text{м/с}$};

\node[flowstart,below=of answer] (finish)
{Готово};

\draw[flowarrow]
(start) -- (data);

\draw[flowarrow]
(data) -- (si);

\draw[flowarrow]
(si) --
node[flowlabel,right]{да}
(sum);

\draw[flowarrow]
(si.west)
-|
node[flowlabel,near start,above]{нет}
(convert.north);

\draw[flowarrow]
(convert.east) -- (sum.west);

\draw[flowarrow]
(sum) -- (formula);

\draw[flowarrow]
(formula) -- (calc);

\draw[flowarrow]
(calc) -- (answer);

\draw[flowarrow]
(answer) -- (finish);

\end{tikzpicture}

\end{center}

\clearpage

% =========================================================
% ТИПИЧНЫЕ ОШИБКИ
% =========================================================

\section*{2. Типичные ошибки}

\begin{itemize}

\item
Сложить скорости отдельных участков и разделить результат на их количество.

\item
Использовать путь только одного участка вместо всего пройденного пути.

\item
Использовать время только одного участка вместо полного времени движения.

\item
Перевести часть величин в СИ, а часть оставить в километрах или минутах.

\item
Перепутать переводы
\[
1\unit{мин}=60\unit{с}
\]
и
\[
1\unit{ч}=3600\unit{с}.
\]

\item
Записать для
$4\unit{мин}$
значение
$200\unit{с}$.
Правильно:
\[
4\cdot60=240\unit{с}.
\]

\item
Потерять один из участков при сложении пути или времени.

\item
Смешать данные разных участков из-за отсутствия индексов.

\item
Сразу подставить числа, пропустив общую формулу.

\item
Оставить окончательный ответ без единицы измерения.

\end{itemize}

\subsection*{Быстрая самопроверка}

Перед сдачей решения задайте себе вопросы:

\begin{enumerate}

\item
Все ли участки пути учтены?

\item
Все ли интервалы времени учтены?

\item
Соответствуют ли друг другу индексы пути и времени?

\item
Все ли величины переведены в совместимые единицы?

\item
Записана ли общая формула?

\item
Есть ли в конце единица измерения?

\end{enumerate}

% =========================================================
% ТРЕНИРОВОЧНЫЙ БЛОК
% =========================================================

\clearpage

\section*{3. Тренировочный блок --- Путь А}

\textbf{Задание 1.}

Велосипедист прошёл три участка:

\[
s_1=120\unit{м},
\qquad
s_2=180\unit{м},
\qquad
s_3=100\unit{м}.
\]

Время движения на участках составило:

\[
t_1=40\unit{с},
\qquad
t_2=60\unit{с},
\qquad
t_3=50\unit{с}.
\]

Найдите среднюю путевую скорость велосипедиста.

\vspace{2em}

\textbf{Задание 2.}

Турист прошёл

\[
0{,}9\unit{км}
\]

за

\[
10\unit{мин},
\]

затем

\[
600\unit{м}
\]

за

\[
5\unit{мин},
\]

после чего прошёл ещё

\[
1{,}5\unit{км}
\]

за

\[
15\unit{мин}.
\]

Найдите среднюю путевую скорость туриста в $\text{м/с}$.

\vspace{2em}

\textbf{Задание 3.}

Лыжник прошёл первый участок длиной

\[
1{,}8\unit{км}
\]

за

\[
20\unit{мин},
\]

второй участок длиной

\[
2{,}4\unit{км}
\]

за

\[
25\unit{мин},
\]

а третий участок длиной

\[
3{,}0\unit{км}
\]

за

\[
35\unit{мин}.
\]

Найдите среднюю путевую скорость лыжника в $\text{м/с}$.

% =========================================================
% ОТВЕТЫ
% =========================================================

\clearpage

\section*{4. Ответы}

\renewcommand{\arraystretch}{1.35}

\begin{table}[h]
\centering

\begin{tabularx}{\linewidth}{@{}cX@{}}

\toprule

\textbf{Задание}
&
\textbf{Ответ}
\\

\midrule

1
&
$\displaystyle
\frac{8}{3}\unit{м/с}
\approx
2{,}67\unit{м/с}$
\\

2
&
$\displaystyle
\frac{5}{3}\unit{м/с}
\approx
1{,}67\unit{м/с}$
\\

3
&
$1{,}5\unit{м/с}$
\\

\bottomrule

\end{tabularx}

\end{table}

\end{document}